% Chapter 29: Time-Varying Treatment and Confounding
% A First Course in Causal Inference - Peng Ding
% Rewritten in Stein motivation-first style

\chapter{时间变动的处理变量与混淆}\label{ch:29}

\section{从单时间点到多时间点：一个自然的推广}\label{sec:29-1}

在前面各章中，我们考虑的都是\textbf{单一时间点}的处理变量——每个单元在某个特定时刻接受或不接受治疗。然而，现实世界中的治疗往往是\textbf{时间变动}的：患者在不同时间点接受不同强度的治疗；政策在不同阶段被动态调整；个人在不同年龄做出不同的选择。

\textbf{为什么时间变动会带来新的挑战？}

考虑一个两时间点的例子：$Z_1$ 是初始治疗（如是否服用阿司匹林），$X_1$ 是治疗后的中间结果（如血压变化），$Z_2$ 是后续治疗（如是否接受手术），$Y$ 是最终结局（如是否在5年内死亡）。在这个序列中：

\begin{itemize}
  \item $Z_1$ 可能影响 $X_1$（治疗改变中间结果）
  \item $X_1$ 可能同时影响 $Z_2$ 和 $Y$（中间结果既是后续治疗的原因，也是结局的原因）
  \item 这时，$X_1$ 同时扮演着\textbf{时间变动协变量}和\textbf{时间变动混淆变量}的角色
\end{itemize}

在单一时间点框架下，我们只需条件于基线协变量 $X_0$ 即可阻断所有混淆。但在时间变动的情况下，$X_1$ 既是治疗 $Z_2$ 的"前置原因"，也是结局 $Y$ 的"前置原因"——它可能本身就是由 $Z_1$ 产生的，从而部分介导了 $Z_1$ 对 $Y$ 的因果效应。这种复杂的时序结构使得简单的条件独立性假设不再够用。

这就是本章要解决的核心问题：\textbf{当处理变量随时间变动时，如何识别和估计因果效应？}

为简化记号，本章主要采用两时间点的设置来阐述核心思想。推广到更多时间点在概念上直接，但有限样本下会产生技术复杂性——组合数随时间点数量指数增长。我们将讨论这些复杂性，并将一般结果留给习题。

\section{基本设置与序列可忽略性}\label{sec:29-2}

\subsection{时序结构}

两时间点的变量时序如下：
\[
X_0 \to Z_1 \to X_1 \to Z_2 \to Y
\]

其中：
\begin{itemize}
  \item $X_0$：基线协变量（治疗前的特征）
  \item $Z_1$：时间点1的处理变量
  \item $X_1$：时间点1与时间点2之间的时变协变量
  \item $Z_2$：时间点2的处理变量
  \item $Y$：最终结局
\end{itemize}

对于二元处理 $(Z_1, Z_2)$，每个单元有四个潜在结局：
\[
Y(z_1, z_2), \quad z_1, z_2 \in \{0, 1\}
\]

观测到的结局等于：
\[
Y = Y(Z_1, Z_2) = \sum_{z_1=0}^1 \sum_{z_2=0}^1 \mathbf{1}(Z_1 = z_1) \mathbf{1}(Z_2 = z_2) Y(z_1, z_2)
\]

\subsection{序列可忽略性假设}

在单一时间点的情况下，$Z \perp Y(z) \mid X$（处理在给定协变量条件下随机化）是我们识别因果效应的关键假设。对于时间变动的情况，这个假设需要\textbf{序列化}——在每个时间点，处理不仅依赖于当前的协变量历史，还依赖于整个历史。

\begin{Assumption}[序列可忽略性]\label{assum:29.1}
\textbf{(1)} $Z_1$ 在给定 $X_0$ 条件下随机化：
\[
Z_1 \perp Y(z_1, z_2) \mid X_0, \quad \forall z_1, z_2 \in \{0, 1\}
\]

\textbf{(2)} $Z_2$ 在给定 $(Z_1, X_1, X_0)$ 条件下随机化：
\[
Z_2 \perp Y(z_1, z_2) \mid (Z_1, X_1, X_0), \quad \forall z_1, z_2 \in \{0, 1\}
\]
\end{Assumption}

\textbf{理解 Assumption \ref{assum:29.1}}：这个假设包含两部分递进的条件。第一部分说明 $Z_1$ 仅被基线协变量 $X_0$ 混淆——一旦我们条件于 $X_0$，$Z_1$ 就等同于随机化实验。第二部分则说明 $Z_2$ 被 $(Z_1, X_1, X_0)$ 混淆——我们不仅需要控制基线协变量 $X_0$，还需要控制由 $Z_1$ 产生的时变协变量 $X_1$。

\textbf{一个重要的微妙之处}：序列可忽略性\textbf{允许} $X_1$ 与 $Y$ 之间存在未测量的混淆 $U$（如图 \ref{fig:29.2} 所示）。这是因为 $X_1$ 本身可能是由 $Z_1$ 产生的，它部分介导了 $Z_1$ 对 $Y$ 的效应——这并非混淆，而是\textbf{中介}。然而，如果 $U$ 同时影响 $X_1$ 和 $Y$（且不通过 $Z_1, Z_2$），则即使满足序列可忽略性，识别仍然会受到挑战。

\begin{figure}[htbp]
\centering
\includegraphics[width=0.6\textwidth]{images/459dbc58690d9628af6d939aa356592d4e8b8cb35fc3ee47e5565a443dbc8cbb.jpg}
\caption{Assumption \ref{assum:29.1} 成立且 $X_1$ 与 $Y$ 之间无未测量混淆 $U$ 的因果图。}
\label{fig:29.1}
\end{figure}

\begin{figure}[htbp]
\centering
\includegraphics[width=0.6\textwidth]{images/fc451469204f8bd2d7424c5f6866f5dbc15053027ff863c8de6b79850192c746.jpg}
\caption{Assumption \ref{assum:29.1} 成立但 $X_1$ 与 $Y$ 之间存在未测量混淆 $U$ 的因果图。序列可忽略性仅排除了 $(Z_1, Z_2)$ 与 $Y$ 之间的混淆，但允许 $X_1$ 与 $Y$ 之间存在未测量关联。}
\label{fig:29.2}
\end{figure}

\section{g公式与结果建模}\label{sec:29-3}

在单一时间点情况下，基于结果的识别公式为：
\[
E\{Y(z)\} = E\{E(Y \mid Z = z, X)\}
\]

对于离散 $X$，这化为：
\[
E\{Y(z)\} = \sum_x E(Y \mid Z = z, X = x) \Pr(X = x)
\]

对于连续 $X$，这化为：
\[
E\{Y(z)\} = \int E(Y \mid Z = z, X = x) f(x) \, dx
\]

接下来的定理将这个公式推广到两时间点的设置。

\begin{Theorem}[g公式]\label{thrm:29.1}
在 Assumption \ref{assum:29.1} 下，
\[
E\{Y(z_1, z_2)\} = E\Big[ E\{E(Y \mid z_2, z_1, X_1, X_0) \mid z_1, X_0\} \Big]
\label{eq:29.1}
\tag{29.1}
\]
\end{Theorem}

\textbf{定理 \ref{thrm:29.1} 的直观理解}：等式右边是从内到外嵌套的条件期望。最内层 $E(Y \mid z_2, z_1, X_1, X_0)$ 是在固定历史 $(z_1, X_0)$、中间协变量 $X_1$ 和第二时刻治疗 $z_2$ 的条件下结局的期望。中间层 $E\{ \cdot \mid z_1, X_0\}$ 对 $X_1$ 求平均——但注意，这是\textbf{逆概率加权}的平均，因为 $\Pr(X_1 \mid z_1, X_0)$ 与 $Z_1 = z_1$ 相关。最外层 $E[\cdot]$ 对 $X_0$ 求边际平均。

对于离散 $X_0, X_1$，\ref{eq:29.1} 化为：
\[
E\{Y(z_1, z_2)\} = \sum_{x_0} \sum_{x_1} E(Y \mid z_2, z_1, x_1, x_0) \Pr(x_1 \mid z_1, x_0) \Pr(x_0)
\label{eq:29.2}
\tag{29.2}
\]

对于连续 $X_0, X_1$，\ref{eq:29.1} 化为：
\[
E\{Y(z_1, z_2)\} = \iint E(Y \mid z_2, z_1, x_1, x_0) f(x_1 \mid z_1, x_0) f(x_0) \, dx_1 \, dx_0
\label{eq:29.3}
\tag{29.3}
\]

\textbf{与全概率公式的比较}：根据观测数据的全概率公式，
\[
\begin{aligned}
E(Y) = & \sum_{x_0} \sum_{z_1} \sum_{x_1} \sum_{z_2} E(Y \mid z_2, z_1, x_1, x_0) \\
& \Pr(z_2 \mid z_1, x_1, x_0) \Pr(x_1 \mid z_1, x_0) \Pr(z_1 \mid x_0) \Pr(x_0)
\end{aligned}
\label{eq:29.4}
\tag{29.4}
\]

对比 \ref{eq:29.3} 与 \ref{eq:29.4}，关键区别在于：\ref{eq:29.3} 中\textbf{没有} $\Pr(z_2 \mid \cdots)$ 和 $\Pr(z_1 \mid \cdots)$ 项——这正是因为潜在结局 $Y(z_1, z_2)$ 的定义是\textbf{固定} $Z_1 = z_1$ 和 $Z_2 = z_2$，而非观测到的分配概率。

\ref{eq:29.2} 和 \ref{eq:29.3} 被称为 \textbf{g公式（g-formula）}，由 Robins 命名。这个名称源于它的一般化形式——当处理变量随时间变动时，这个公式提供了一种递推的识别策略。

\textbf{定理 \ref{thrm:29.1} 的证明}。\footnote{详见附录 \cref{sec:derivation-thm-29.1}。}

\subsection{基于结果建模的插件估计}\label{sec:29-3-1}

g公式 \ref{eq:29.2} 和 \ref{eq:29.3} 提示我们：为了估计潜在结局的均值，我们需要对 $E(Y \mid z_2, z_1, x_1, x_0)$、$\Pr(x_1 \mid z_1, x_0)$ 和 $\Pr(x_0)$ 建模。有了这些拟合模型，我们可以将它们"插入"g公式得到估计。

在某些特殊函数形式下，这个任务可以简化。下面的例29.1给出线性结果模型下的结果。

\begin{Example}[例29.1 — 线性结果模型]\label{ex:29.1}
假设线性结果模型：
\[
E(Y \mid z_2, z_1, x_1, x_0) = \beta_0 + \beta_1 z_2 + \beta_2 z_1 + \beta_3 x_1 + \beta_4 x_0
\tag{29.5}
\]

可以验证：
\[
\begin{aligned}
E\{Y(z_1, z_2)\} = & \sum_{x_0} \sum_{x_1} (\beta_0 + \beta_1 z_2 + \beta_2 z_1 + \beta_3 x_1 + \beta_4 x_0) \\
& \Pr(x_1 \mid z_1, x_0) \Pr(x_0) \\
= & \beta_0 + \beta_1 z_2 + \beta_2 z_1 + \beta_3 \sum_{x_0} E(X_1 \mid z_1, x_0) \Pr(x_0) + \beta_4 E(X_0) \\
= & \beta_0 + \beta_1 z_2 + \beta_2 z_1 + \beta_3 E\{X_1(z_1)\} + \beta_4 E(X_0)
\end{aligned}
\tag{29.6}
\]

其中我们定义了 $X_1$ 在治疗 $Z_1 = z_1$ 下的潜在结局：
\[
E\{X_1(z_1)\} = \sum_{x_0} E(X_1 \mid z_1, x_0) \Pr(x_0)
\tag{29.7}
\]

定义 $\tau_{Z_1 \to X_1} = E\{X_1(1) - X_1(0)\}$，则：
\[
E\{Y(1, 0) - Y(0, 0)\} = \beta_2 + \beta_3 \tau_{Z_1 \to X_1}
\label{eq:29.8}
\tag{29.8}
\]
\[
E\{Y(0, 1) - Y(0, 0)\} = \beta_1
\label{eq:29.9}
\tag{29.9}
\]
\[
E\{Y(1, 1) - Y(0, 0)\} = \beta_1 + \beta_2 + \beta_3 \tau_{Z_1 \to X_1}
\label{eq:29.10}
\tag{29.10}
\]

\textbf{直观解释}（参见图 \ref{fig:29.3}）：\ref{eq:29.8} 中 $Z_1$ 的效应等于路径 $Z_1 \to Y$（系数 $\beta_2$）和路径 $Z_1 \to X_1 \to Y$（系数 $\beta_3 \tau_{Z_1 \to X_1}$）之和；\ref{eq:29.9} 中 $Z_2$ 的效应等于路径 $Z_2 \to Y$（系数 $\beta_1$）；\ref{eq:29.10} 中 $(Z_1, Z_2)$ 的总效应等于三条路径系数之和。

\begin{figure}[htbp]
\centering
\includegraphics[width=0.5\textwidth]{images/713d1ba2f50c0527a6a7759f58f6e5942115a4ac9b2746ad7874cc5347518c86.jpg}
\caption{带路径系数的线性因果图，条件于基线协变量 $X_0$。}
\label{fig:29.3}
\end{figure}
\end{Example}

\textbf{g-null 悖论}。Robins 和 Wasserman（1997）指出了基于结果建模的插件估计的一个令人惊讶的缺陷。他们证明：如果在这个策略中出现模型错误设定，数据分析者可能\textbf{错误地拒绝}零因果效应 null 假设，即使在真实数据生成过程中 $(Z_1, Z_2)$ 对 $Y$ 的因果效应确实为零。他们称之为 \textbf{g-null 悖论（g-null paradox）}。令人惊讶的是，即使在例29.1的简单线性结果模型中，g-null 悖论也可能出现（见习题29.1）。

\subsection{基于结果建模的递推估计}\label{sec:29-3-2}

29.3.1节的插件估计涉及对时变协变量 $X_1$ 建模，并导致令人不快的 g-null 悖论——这不是一个理想的方法。

递推期望公式 \ref{eq:29.1} 启发了一种更简单的估计方法。从最内层条件期望开始，记：
\[
\tilde{Y}_2(z_1, z_2) = E(Y \mid Z_2 = z_2, Z_1 = z_1, X_1, X_0)
\]

我们可以用 $(Z_2 = z_2, Z_1 = z_1)$ 子集数据拟合 $Y$ 对 $(X_1, X_0)$ 的模型，得到所有单元的拟合值 $\hat{Y}_{2i}(z_1, z_2)$。然后进入外层条件期望，记：
\[
\tilde{Y}_1(z_1, z_2) = E\{\tilde{Y}_2(z_1, z_2) \mid Z_1 = z_1, X_0\}
\]

我们可以用 $Z_1 = z_1$ 子集数据拟合 $\hat{Y}_2(z_1, z_2)$ 对 $X_0$ 的模型，得到所有单元的拟合值 $\hat{Y}_{1i}(z_1, z_2)$。最终的 $E\{Y(z_1, z_2)\}$ 估计量为：
\[
\hat{E}\{Y(z_1, z_2)\} = n^{-1} \sum_{i=1}^n \hat{Y}_{1i}(z_1, z_2)
\tag{29.11}
\]

这种递推估计\textbf{不涉及}对 $X_1$ 建模，因此避免了 g-null 悖论（见习题29.2）。然而，基于递推回归的估计量实现起来并不容易，因为它涉及建模与因果图自然结构不对应的变量——例如 $\tilde{Y}_2(z_1, z_2)$。

\section{逆概率加权}\label{sec:29-4}

在单一时间点情况下，IPW 识别公式为：
\[
E\{Y(z)\} = E\left\{ \frac{\mathbf{1}(Z = z) Y}{\Pr(Z = z \mid X)} \right\}
\]

接下来的结果将其推广到两时间点设置。定义：
\[
e(z_1, X_0) = \Pr(Z_1 = z_1 \mid X_0)
\tag{29.12}
\]
\[
e(z_2, Z_1, X_1, X_0) = \Pr(Z_2 = z_2 \mid Z_1, X_1, X_0)
\tag{29.13}
\]

为时间点1和时间点2的倾向得分。

\begin{Theorem}[IPW 识别公式]\label{thrm:29.2}
在 Assumption \ref{assum:29.1} 下，
\[
E\{Y(z_1, z_2)\} = E\left\{ \frac{\mathbf{1}(Z_1 = z_1) \mathbf{1}(Z_2 = z_2) Y}{e(z_1, X_0) \, e(z_2, Z_1, X_1, X_0)} \right\}
\label{eq:29.14}
\tag{29.14}
\]
\end{Theorem}

\textbf{定理 \ref{thrm:29.2} 揭示的重叠假设}：
\[
0 < e(z_1, X_0) < 1, \quad 0 < e(z_2, Z_1, X_1, X_0) < 1
\label{eq:29.15}
\tag{29.15}
\]

对所有 $z_1, z_2$。如果某些倾向得分为 0 或 1，识别公式 \ref{eq:29.14} 会爆炸到无穷。

\textbf{直观理解}：这个公式本质上是单一时间点 IPW 公式的两次应用——首先在 $Z_2$ 上"调整"，然后在 $Z_1$ 上"调整"。分子中的指示函数 $\mathbf{1}(Z_1 = z_1) \mathbf{1}(Z_2 = z_2)$ 选择出接受特定治疗序列的单元，分母中的倾向得分对此进行加权，使每个单元在分析中代表与其倾向得分成反比的单元数。

\textbf{估计策略}。基于 IPW 的估计非常简单，只需对两个二元治疗指标建模。首先，我们可以拟合 $Z_1$ 对 $X_0$ 的模型，得到所有单元的拟合值 $\hat{e}_1(z_1, X_{0i})$；拟合 $Z_2$ 对 $(Z_1, X_1, X_0)$ 的模型，得到所有单元的拟合值 $\hat{e}_2(z_2, Z_{1i}, X_{1i}, X_{0i})$。然后得到以下 IPW 估计量（Horvitz-Thompson 型）：
\[
\hat{E}^{\mathrm{ht}}\{Y(z_1, z_2)\} = n^{-1} \sum_{i=1}^n \frac{\mathbf{1}(Z_{1i} = z_1) \mathbf{1}(Z_{2i} = z_2) Y_i}{\hat{e}_1(z_1, X_{0i}) \hat{e}_2(z_2, Z_{1i}, X_{1i}, X_{0i})}
\tag{29.16}
\]

与第11章的讨论类似，HT 估计量对结局的位移不具有不变性，且在有限样本下存在不稳定性。修正后的 Hajek 型估计量为 $\hat{E}^{\mathrm{haj}}\{Y(z_1, z_2)\} = \hat{E}^{\mathrm{ht}}\{Y(z_1, z_2)\} / \hat{1}^{\mathrm{ht}}(z_1, z_2)$，其中：
\[
\hat{1}^{\mathrm{ht}}(z_1, z_2) = n^{-1} \sum_{i=1}^n \frac{\mathbf{1}(Z_{1i} = z_1) \mathbf{1}(Z_{2i} = z_2)}{\hat{e}_1(z_1, X_{0i}) \hat{e}_2(z_2, Z_{1i}, X_{1i}, X_{0i})}
\tag{29.17}
\]

\section{多时间点情形}\label{sec:29-5}

将29.2和29.3节的估计策略推广到多时间点并非易事。即使对于二元处理和 $K$ 个时间点，处理组合的数量也随 $K$ 指数增长（例如 $2^5 = 32$，$2^{10} = 1024$）。因此，29.2和29.3节的结果回归和 IPW 估计量在有限样本下是不可行的——因为它们要求每个处理组合水平组合都有足够的数据。

为应对这一挑战，Robins 提出了两种强大的估计框架：\textbf{边际结构模型（Marginal Structural Model, MSM）}和\textbf{结构嵌套模型（Structural Nested Model, SNM）}。

\subsection{边际结构模型}\label{sec:29-5-1}

一种强有力的方法基于\textbf{边际结构模型（MSM）}（Robins et al., 2000; Hernán et al., 2000）。为简化记号，我们仅展示 $K=2$ 的 MSM，尽管其主要用途在于一般情形。

\begin{Definition}[MSM — 边际结构模型]\label{def:29.1}
$Y(z_1, z_2)$ 的边际均值为：
\[
E\{Y(z_1, z_2)\} = f(z_1, z_2; \beta)
\tag{29.18}
\]
\end{Definition}

一个典型例子是 $E\{Y(z_1, z_2)\} = \beta_0 + \beta_1 z_1 + \beta_2 z_2$。我们也可以将基线协变量纳入模型。

\begin{Definition}[带基线协变量的 MSM]\label{def:29.2}
给定 $X_0$ 条件下 $Y(z_1, z_2)$ 的均值为：
\[
E\{Y(z_1, z_2) \mid X_0\} = f(z_1, z_2, X_0; \beta)
\tag{29.19}
\]
\end{Definition}

一个典型例子是：
\[
E\{Y(z_1, z_2) \mid X_0\} = \beta_0 + \beta_1 z_1 + \beta_2 z_2 + \beta_3^{\mathsf{T}} X_0
\label{eq:29.20}
\tag{29.20}
\]

如果我们观测到所有潜在结局，可以通过以下最小化问题求解 $\beta$：
\[
\beta = \arg\min_b \sum_{z_2} \sum_{z_1} E\left\{ Y(z_1, z_2) - f(z_1, z_2, X_0; b) \right\}^2
\label{eq:29.21}
\tag{29.21}
\]

在序列可忽略性下，我们可以通过仅涉及观测变量的最小化问题求解 $\beta$。

\begin{Theorem}[MSM 下的 IPW]\label{thrm:29.3}
在 Assumption \ref{assum:29.1} 和 Definition \ref{def:29.2} 下，\ref{eq:29.21} 中的 $\beta$ 等于：
\[
\beta = \arg\min_b \sum_{z_2} \sum_{z_1} E\left[ \frac{\mathbf{1}(Z_1 = z_1) \mathbf{1}(Z_2 = z_2)}{e(z_1, X_0) \, e(z_2, Z_1, X_1, X_0)} \{ Y - f(z_1, z_2, X_0; b) \}^2 \right]
\tag{29.22}
\]
\end{Theorem}

\textbf{估计策略}。定理 \ref{thrm:29.3} 暗示了一种基于加权回归的简单估计策略。例如，在 \ref{eq:29.20} 下，我们可以拟合 $Y_i$ 对 $(1, Z_{1i}, Z_{2i}, X_{0i})$ 的加权最小二乘（WLS）回归，权重为 $\hat{e}_1^{-1}(Z_{1i}, X_{0i}) \hat{e}_{2i}^{-1}(Z_{2i}, Z_{1i}, X_{1i}, X_{0i})$。

\textbf{与结果回归的比较}：MSM 的 IPW 方法与结果回归方法的关键区别在于：结果回归直接建模 $E(Y \mid Z, X)$，而 MSM 的 IPW 方法建模 $E\{Y(z) \mid X_0\}$ 并通过逆概率加权来处理时变混淆。MSM 的优势在于它允许我们指定\textbf{边际}（而非条件）结构，这有时更加自然。

\textbf{高维协变量下的 MSM}。当 $X_0$ 维数很高时，\ref{eq:29.20} 中的线性模型 $\beta_3^{\mathsf{T}} X_0$ 面临维数灾难。此时，稀疏性假设变得关键——即 $\beta_3$ 中只有少数元素非零。LASSO（Tibshirani, 1996）等正则化方法可以用于在高维下估计 $\beta$，但需要额外的理论保证来确保因果推断的渐近有效性（见 Bloniarz et al., 2016; Wager et al., 2016）。

\subsection{结构嵌套模型}\label{sec:29-5-2}

IPW 的一个关键问题是：如果重叠假设 \ref{eq:29.15} 被违背，IPW 不可用。为解决这一挑战，Robins 提出了结构嵌套模型。这里我们仅回顾两时间点的版本。

\begin{Definition}[结构嵌套模型（SNM）]\label{def:29.3}
时间点1的条件效应为：
\[
E\{Y(z_1, 0) - Y(0, 0) \mid Z_1 = z_1, X_0\} = g_1(z_1, X_0; \beta)
\tag{29.23}
\]

时间点2的条件效应为：
\[
E\{Y(z_1, z_2) - Y(z_1, 0) \mid Z_2 = z_2, Z_1 = z_1, X_1, X_0\} = g_2(z_2, z_1, X_1, X_0; \beta)
\tag{29.24}
\]
\end{Definition}

在 Definition \ref{def:29.3} 中，有两个逻辑约束：
\[
g_1(0, X_0; \beta) = 0, \quad g_2(0, z_1, X_1, X_0; \beta) = 0, \quad \forall z_1
\tag{29.25}
\]

两个典型选择如下。

\begin{Example}[例29.2 — 加法齐次 SNM]\label{ex:29.2}
假设：
\[
\begin{cases}
g_1(z_1, X_0; \beta) = \beta_1 z_1 \\
g_2(z_2, z_1, X_1, X_0; \beta) = (\beta_2 + \beta_3 z_1) z_2
\end{cases}
\tag{29.26}
\]
\end{Example}

\begin{Example}[例29.3 — 协变量调整 SNM]\label{ex:29.3}
假设：
\[
\begin{cases}
g_1(z_1, X_0; \beta) = (\beta_1 + \beta_2^{\mathsf{T}} X_0) z_1 \\
g_2(z_2, z_1, X_1, X_0; \beta) = (\beta_3 + \beta_4 z_1 + \beta_5^{\mathsf{T}} X_1) z_2
\end{cases}
\tag{29.27}
\]
\end{Example}

\textbf{MSM 与 SNM 的比较}。结构嵌套模型允许同时调整基线协变量和时变协变量，而边际结构模型仅允许调整基线协变量。Definition \ref{def:29.3} 下的估计更为复杂，策略基于估计方程。

\textbf{估计方程}。首先引入两个重要构建块。定义：
\[
U_2(\beta) = Y - g_2(Z_2, Z_1, X_1, X_0; \beta)
\tag{29.28}
\]
\[
U_1(\beta) = Y - g_2(Z_2, Z_1, X_1, X_0; \beta) - g_1(Z_1, X_0; \beta)
\tag{29.29}
\]

它们不能直接从数据计算，因为它们依赖于参数的真实值。在真值处，它们具有以下性质。

\begin{Lemma}[引理 29.1]\label{lemm:29.1}
在 Assumption \ref{assum:29.1} 和 Definition \ref{def:29.3} 下，
\[
E\{U_2(\beta) \mid Z_2, Z_1, X_1, X_0\} = E\{U_2(\beta) \mid Z_1, X_1, X_0\} = E\{Y(Z_1, 0) \mid Z_1, X_1, X_0\}
\tag{29.30}
\]
且
\[
E\{U_1(\beta) \mid Z_1, X_0\} = E\{U_1(\beta) \mid X_0\} = E\{Y(0, 0) \mid X_0\}
\tag{29.31}
\]
\end{Lemma}

\textbf{引理 \ref{lemm:29.1} 的直观理解}：$U_1(\beta)$ 扮演着"接受任何治疗前的对照潜在结局"的角色，而 $U_2(\beta)$ 扮演着"接受时间点1治疗后、时间点2前的对照潜在结局"的角色。

\begin{Theorem}[定理 29.4 — SNM 估计方程]\label{thrm:29.4}
在 Assumption \ref{assum:29.1} 和 Definition \ref{def:29.3} 下，
\[
E\left[ h_2(Z_1, X_1, X_0) \{ Z_2 - e(1, Z_1, X_1, X_0) \} U_2(\beta) \right] = 0
\tag{29.32}
\]
且
\[
E\left[ h_1(X_0) \{ Z_1 - e(1, X_0) \} U_1(\beta) \right] = 0
\tag{29.33}
\]

对任何函数 $h_1$ 和 $h_2$（只要矩存在）。
\end{Theorem}

使用定理 \ref{thrm:29.4} 时，我们必须指定 $h_1$ 和 $h_2$ 以确保有足够的方程来求解 $\beta$。例29.4回顾例29.2。

\begin{Example}[例29.4 — 例29.2下的估计方程]\label{ex:29.4}
在例29.2下，我们可以选择 $h_1 = 1$ 和 $h_2 = (1, Z_1)$ 得到：
\[
\begin{aligned}
E[\{Z_2 - e(1, Z_1, X_1, X_0)\} \{Y - (\beta_2 + \beta_3 Z_1) Z_2\}] & = 0 \\
E[Z_1 \{Z_2 - e(1, Z_1, X_1, X_0)\} \{Y - (\beta_2 + \beta_3 Z_1) Z_2\}] & = 0 \\
E[\{Z_1 - e(1, X_0)\} \{Y - (\beta_2 + \beta_3 Z_1) Z_2 - \beta_1 Z_1\}] & = 0
\end{aligned}
\tag{29.34}
\]

然后我们可以从上述线性方程组求解 $\beta$（见习题29.6）。
\end{Example}

\textbf{更高效的估计}。一个自然的问题是：替代性的 $(h_1, h_2)$ 选择是否能导致更高效的估计量？答案是肯定的。例如，我们可以选择许多 $(h_1, h_2)$ 并使用广义矩方法（GMM; Hansen, 1982）。技术细节超出本书范围。Naimi et al.（2017）和 Vansteelandt and Joffe（2014）提供了结构嵌套模型的教程。

\section{习题}\label{sec:29-homework}

\begin{HomeworkProblem}[g-null paradox]\label{exr:29-1}


Consider the simple causal diagram in Figure \ref{fig::si-g-null-paradox} without pre-treatment covariates $X_0$ and without the arrows from $(Z_1, Z_2)$ to $Y$. So the effect of $(Z_1,Z_2)$ on $Y$ is zero.

\begin{figure}
\bigskip
\[
 \xymatrix{
 Z_1 \ar[r]   & X_1   \ar[r]  & Z_2  & Y \\
                                             &                                             &U\ar[lu]\ar[ru]
 }
 \]
 \caption{With unmeasured confounding between $X_1$ and $Y$. The causal diagram ignores the pre-treatment covariates $X_0$.}\label{fig::si-g-null-paradox}
\end{figure}


Revisit Example \ref{example::linear-si}. Show that the expectation $E\{ Y(z_1, z_2) \} $ does not depend on $(z_1, z_2)$ if
\[
\beta_1  = \beta_2 = 0  \text{ and }   \beta_3 = 0
\]
or
\[
\beta_1 = \beta_2 = 0 \text{ and } E\{ X_1(z_1)\}  \text{ does not depend on } z_1.
\]
holds.


Remark:
However, $\beta_3 = 0$ in the first condition rules out the dependence of $Y$ on $X_1$, contradicting the existence of unmeasured confounder $U$ between $X_1$ and $Y$; the independence of $E\{ X_1(z_1)\} $ on $z_1$ rules out the dependence of $X_1$ on $Z_1$, contradicting with the existence of the arrow from $Z_1$ on $X_1$. That is, if there is an unmeasured confounder $U$ between $X_1$ and $Y$ and there is an arrow from $Z_1$ on $X_1$, then the formula of $E\{ Y(z_1, z_2) \} $ in Example \ref{example::linear-si} must depend on $(z_1, z_2)$, which leads to a contradiction with the absence of arrows from $(Z_1, Z_2)$ to $Y$.


\end{HomeworkProblem}

\begin{HomeworkProblem}[Recursive estimation under the null model]\label{exr:29-2}


Consider the recursive estimation method in \ref{section::recursive} under the causal diagram in Problem \ref{problem::g-null-paradox}.
Show that based on linear models, the estimator converges to 0.



\end{HomeworkProblem}

\begin{HomeworkProblem}[IPW under MSM]\label{exr:29-3}


Prove Theorem \ref{thm::si-msm-ipw}.


\end{HomeworkProblem}

\begin{HomeworkProblem}[A nonlinear example of Definition \ref{def::si-msm-with-x0}]\label{exr:29-4}


Another leading example of Definition \ref{def::si-msm-with-x0} is
\begin{equation}\label{eq::logistic-msm}
\text{logit} \left[ \pr\{  Y(z_1, z_2)  =1  \mid X_0\}  \right] = \beta_0 + \beta_1 z_1 + \beta_2 z_2 + \beta_3\tran X_0 .
\end{equation}
If we observe all potential outcomes, we can solve $\beta$ by minimizing the expectation of the negative log-likelihood function (see Chapter \ref{sec::mle-logistic} for a simpler version):
\begin{equation}
\beta = \arg\min_b   \sum_{z_2}\sum_{z_1}  E\left\{   \log (1 + e^{\ell })
- Y(z_1, z_2)  \ell  \right\}
\end{equation}
where $\ell = \beta_0 + \beta_1 z_1 + \beta_2 z_2 + \beta_3\tran X_0.$
Under sequential ignorability, we can solve $\beta$ from the following minimization problem that only involves observables.


\begin{Theorem}
[IPW under MSM]
\label{thm::si-msm-ipw-logistic}
Under Assumption \ref{assume::sr-time-varying} and Definition \ref{def::si-msm-with-x0}, the $\beta$ in \eqref{eq::beta-logistic} equals
\[
\beta = \arg\min_b \sum_{z_2}\sum_{z_1}  E\left[   \frac{  1(Z_1=z_1) 1(Z_2=z_2)   }{   e(z_1, X_0) e(z_2, Z_1, X_1, X_0)     }
\left\{   \log (1 + e^{\ell })
- Y(z_1, z_2)  \ell  \right\} \right] .
\]
\end{Theorem}

Prove Theorem \ref{thm::si-msm-ipw-logistic}.

Remark:
Theorem \ref{thm::si-msm-ipw-logistic} implies a simple estimation strategy based on weighted regressions. For instance, under \eqref{eq::logistic-msm}, we can fit weighted logistic regression of $Y_i$ on $(1,Z_{1i}, Z_{2i}, X_{0i})$ with weights $\hat{e}_{1}^{-1}(Z_{1i}, X_{0i})  \hat{e}_{2i}^{-1}(Z_{2i}, Z_{i1}, X_{1i}, X_{0i})$.






\end{HomeworkProblem}

\begin{HomeworkProblem}[Structural nested model with a single time point]\label{exr:29-5}



Recall the standard setting of observational studies with IID data drawn from $\{ X,Z, Y(1), Y(0)\}$. Define the propensity score as $e(X) = \pr(Z=1\mid X)$. Assume
\[
Z\ind Y(0)\mid X
\]
and the following structural nested model.


\begin{definition}
[structural nested model with a single time point]
\label{definition::snm-single}
The conditional mean of the individual effect is
\[
E\{  Y(z) - Y(0)\mid Z=z, X \} = g(z,X;\beta).
\]
\end{definition}

In Definition \ref{definition::snm-single}, a logical restriction is $g(0,X;\beta) = 0$. Prove the following results.


\begin{enumerate}
\item
We have
\[
E\{  Y - g(Z,X;\beta) \mid X, Z\} = E\{  Y - g(Z,X;\beta) \mid X\}  = E\{  Y(0) \mid X\}.
\]


\item
We have
\begin{eqnarray}\label{eq::Estimating-equation-snm-1}
E\Big [  h(X)  \{ Z-e(X) \}   \{  Y - g(Z,X;\beta) \}  \Big ] = 0
\end{eqnarray}
for any function $h$, provided that the moment exists.
\end{enumerate}




Remark:  Equation \eqref{eq::Estimating-equation-snm-1} is the basis for parameter estimation.
Consider a special case of Definition \ref{definition::snm-single} with $g(z,X;\beta) = \beta z$. Choose $h(X) = 1$ to obtain
\[
E\{   (Z-e(X))   (Y-  \beta Z) \}  = 0.
\]
Solve for $\beta$ to obtain
\[
\beta = \frac{   E\{ (Z-e(X)) Y \}   }{   E\{ (Z-e(X))  Z\}   }.
\]
Therefore, $\beta$ equals the coefficient of $Z$ in the TSLS of $Y$ on $Z$ with $Z-e(X)$ being the IV for $Z$. With some basic calculations, we can also show that
\[
\beta = \frac{   \cov\{ Z-e(X), Y\}   }{ \cov\{ Z-e(X) \}  }.
\]
Therefore, $\beta$ equals the coefficient of $Z-e(X)$ in the OLS of $Y$ on $Z-e(X)$, which appeared in Chapter \ref{sec::regress-on-pscore} before.

Consider another special case of Definition \ref{definition::snm-single} with $g(z,X;\beta) = (\beta_0 + \beta_1\tran X) z$. Choose $h(X) = (1,X)$ to obtain
\[
E\left\{   \begin{pmatrix}
Z-e(X)\\
(Z-e(X)) X
\end{pmatrix}
(Y-  \beta_0 Z - \beta_1\tran XZ) \right\}  = 0.
\]
That is, $(\beta_0,\beta_1)$ equal the coefficients in the TSLS of $Y$ on $(Z, XZ)$ with $(Z-e(X),  (Z-e(X)) X )$ being the IV for $(Z, XZ)$.



\end{HomeworkProblem}

\begin{HomeworkProblem}[Estimation under Example \ref{example::g-estimation-snm}]\label{exr:29-6}


We can estimate the $\beta$'s by solving the empirical version of the estimating equations in Example \ref{example::g-estimation-snm}. We first estimate the two propensity scores and obtain the centered treatment
\[
\check{Z}_{1i} = Z_{1i} - \hat{e}(1,X_{0i})
\]
 at time point 1 and
 \[
 \check{Z}_{2i} = Z_{2i}  - \hat{e}(1, Z_{1i}, X_{1i}, X_{0i})
 \]
 at time point 2.


Show that we can estimate $\beta_2$ and $\beta_3$ by running TSLS of $Y_i$ on $(Z_{2i}, Z_{1i} Z_{2i} )$ with $( \check{Z}_{2i} ,  Z_{1i}\check{Z}_{2i} )$ as the IV for $(Z_{2i}, Z_{1i} Z_{2i} )$, and then we can estimate $\beta_1$ by running TSLS of $Y_i - (\hat{\beta}_2+\hat{\beta}_3 Z_{1i}) Z_{2i}$ on $Z_{1i}$ with $\check{Z}_{1i} $ as the IV for $Z_{1i}$.




\end{HomeworkProblem}

\begin{HomeworkProblem}[g-formula with a treatment at multiple time points]\label{exr:29-7}



Extend the discussion to the setting with $K$ time points. The temporal ordering (but not the causal diagram) of the variables is
\[
X_0 \rightarrow Z_1 \rightarrow X_1 \rightarrow Z_2 \rightarrow \cdots  \rightarrow  X_{K-1}  \rightarrow Z_K .
\]
Introduce the notation $\overline{Z}_k = (Z_1, \ldots, Z_k)$ and $\overline{X}_k = (X_0, X_1, \ldots, X_k)$ with lower case $\overline{z}_k $ and $\overline{x}_k$  denoting the corresponding  realized values. With $k=0$, we have $\overline{X}_0 = X_0$ and $\overline{Z}_0$ is empty.
Each unit has $2^K$ potential outcomes:
\[
Y(\overline{z}_K) \text{ for all } z_1,\ldots, z_K = 0,1.
\]
Assume sequential ignorability below.

\begin{assumption}
[sequential ignorability at multiple time points]\label{assume::si-multiple}
We have
\[
Z_k \ind Y(\overline{z}_K) \mid (\overline{Z}_{k-1}, \overline{X}_{k-1})
\]
for all $k=1,\ldots, K$ and all $z_1,\ldots, z_K = 0,1.$
\end{assumption}


Prove Theorem \ref{thm::g-formula-multiple} below.


\begin{Theorem}
[g-formula with multiple time points]
\label{thm::g-formula-multiple}
Under Assumption \ref{assume::si-multiple},
\[
E\{ Y(\overline{z}_K)  \} = E\left[ \cdots
E\{   E(Y\mid \overline{z}_K, \overline{X}_{K-1}) \mid \overline{z}_{K-1}, \overline{X}_{K-2} \}
\cdots \mid z_1,  X_0 \right] .
\]
\end{Theorem}


Remark:
In Theorem \ref{thm::g-formula-multiple}, I use the simplified notation ``$\overline{z}_k$'' for  ``$\overline{Z}_k = \overline{z}_k $.''
With discrete $X$, Theorem \ref{thm::g-formula-multiple} reduces to
\begin{eqnarray*}
 E\{ Y(\overline{z}_K)  \}
&=&
\sum_{x_0} \sum_{x_1} \cdots   \sum_{x_{K-1}}
 E(Y\mid \overline{z}_K, \overline{x}_{K-1})   \\
 &&
\qquad \qquad \cdot  \pr(x_{K-1}\mid \overline{z}_{K-1}, \overline{x}_{K-2})
 \cdots
 \pr(x_1\mid z_{1}, x_0)
 \pr(x_0) ;
\end{eqnarray*}
with continuous $X$, Theorem \ref{thm::g-formula-multiple} reduces to
\begin{eqnarray*}
  E\{ Y(\overline{z}_K)  \}
&=&
\int
 E(Y\mid \overline{z}_K, \overline{x}_{K-1})   \\
 && \cdot
 f(x_{K-1}\mid \overline{z}_{K-1}, \overline{x}_{K-2})
 \cdots
 f(x_1\mid z_{1}, x_0)
 f(x_0)
 \d \overline{x}_{K-1}   .
\end{eqnarray*}






\end{HomeworkProblem}

\begin{HomeworkProblem}[IPW with treatments at multiple time points]\label{exr:29-8}

Inherit the setting of Problem \ref{problem::g-formula-multiple-time}. Define the propensity score at $K$ time points as
\begin{eqnarray*}
e(z_1, X_0) &=& \pr(Z_1 = z_1\mid X_0) ,\\
\vdots \\
e(z_k, \overline{Z}_{k-1},  \overline{X}_{k-1}) &=&  \pr(Z_k = z_k \mid \overline{Z}_{k-1},  \overline{X}_{k-1}) ,\\
\vdots \\
e(z_K, \overline{Z}_{K-1},  \overline{X}_{K-1}) &=&  \pr(Z_K = z_K \mid \overline{Z}_{K-1},  \overline{X}_{K-1}) .
\end{eqnarray*}


Prove Theorem \ref{thm::ipw-multiple} below assuming overlap implicitly.


\begin{Theorem}[IPW with multiple time points]\label{thm::ipw-multiple}
Under Assumption \ref{assume::si-multiple},
\[
  E\{ Y(\overline{z}_K)  \}   = E\left\{
  \frac{1(Z_1 = z_1) \cdots  1( Z_K = z_K) Y}{    e(z_1, X_0)  \cdots  e(z_K, \overline{Z}_{K-1},  \overline{X}_{K-1})   }
  \right\}.
\]
\end{Theorem}


Based on Theorem \ref{thm::ipw-multiple}, construct the Horvitz--Thompson and Hajek estimators.


\end{HomeworkProblem}

\begin{HomeworkProblem}[MSM with treatments at multiple time points]\label{exr:29-9}


The number of potential outcomes grows exponentially with $K$. The formulas in Problems \ref{problem::g-formula-multiple-time} and \ref{problem::ipw-multiple-time} are not directly applicable in finite samples. We can impose the following structural assumptions on the potential outcomes.


\begin{definition}
[MSM with multiple time points]\label{definition::msm-multiple}
Assume
\[
  E\{ Y(\overline{z}_K)  \mid X_0\}  = f(\overline{z}_K, X_0 ; \beta  ).
\]
\end{definition}


Two leading examples of Definition \ref{definition::msm-multiple} are
\[
  E\{ Y(\overline{z}_K)  \mid X_0\}   = \beta_0 + \beta_1 \sum_{k=1}^K z_k + \beta_2\tran X_0
  \]
and
\[
 E\{ Y(\overline{z}_K)  \mid X_0\}   = \beta_0 +  \sum_{k=1}^K \beta_k z_k + \beta_{K+1}\tran X_0 .
 \]

If we know all the potential outcomes, we can solve $\beta$ from the following minimization problem:
\[
\beta = \arg \min_b  \sum_{  \overline{z}_K  }  E\{   Y(\overline{z}_K) -  f(\overline{z}_K, X_0 ; \beta  )\}^2.
\]
Theorem \ref{thm::ipw-msm-multiple} below shows that under Assumption \ref{assume::si-multiple}, we can solve $\beta$ from a minimization problem that only involves observables.


\begin{Theorem}[IPW for MSM with multiple time points]\label{thm::ipw-msm-multiple}
Under Assumption \ref{assume::si-multiple},
\[
\beta = \arg \min_b  \sum_{  \overline{z}_K  }    E\left[
  \frac{1(Z_1 = z_1) \cdots 1( Z_K = z_K) }{    e(z_1, X_0)  \cdots  e(z_K, \overline{Z}_{K-1},  \overline{X}_{K-1})    }
  \{ Y -  f(\overline{z}_K, X_0 ; \beta  )\} ^2
  \right] .
\]
\end{Theorem}




\end{HomeworkProblem}

\begin{HomeworkProblem}[Structural nested model with treatments at multiple time points]\label{exr:29-10}


Inherit the setting from Problem \ref{problem::g-formula-multiple-time} and the notation from Problem \ref{problem::ipw-multiple-time}. This problem presents a general structural nested model.

\begin{definition}
[structural nested model with multiple time points]
\label{def::snm-multiple}
The conditional effect at time $k$ is
\[
E\{  Y(\overline{z}_k,0) - Y( \overline{z}_{k-1},0) \mid \overline{z}_k , \overline{X}_{k-1} \} = g_k( \overline{z}_k, \overline{X}_{k-1}; \beta)
\]
for all $\overline{z}_k$ and all $k=1,\ldots, K$.
\end{definition}



In Definition \ref{def::snm-multiple},  a logical restriction is
\[
g_k(0,  \overline{z}_{k-1}, \overline{X}_{k-1}; \beta)    = 0
\]
for all $\overline{z}_{k-1}$ and all $k=1,\ldots, K$.


Define
\[
U_k(\beta) = Y - \sum_{s=1}^k  g_s( \overline{Z}_s, \overline{X}_{s-1}; \beta)
\]
for all $k=1,\ldots, K$. Theorem  \ref{thm::snm-estimating-equation-multiple} below extends Theorem \ref{thm::estimating-equation-snm}.



\begin{Theorem}
\label{thm::snm-estimating-equation-multiple}
Under Assumption \ref{assume::si-multiple} and Definition \ref{def::snm-multiple},  we have
\[
E\left[    h_k( \overline{Z}_{k-1}, \overline{X}_{k-1} ) \{  Z_k - e(1, \overline{Z}_{k-1}, \overline{X}_{k-1}) \}  U_k(\beta)
\right] = 0
\]
for any functions $h_k$ $(k=1,\ldots, K)$, provided that   the moment exists.
\end{Theorem}



Remark: Choosing appropriate $h_k$'s, we can estimate $\beta$ by solving the empirical version of Theorem \ref{thm::snm-estimating-equation-multiple}.



\end{HomeworkProblem}

\begin{HomeworkProblem}[Recommended reading]\label{exr:29-11}

\citet{robins2000marginal} reviewed the MSM. \citet{naimi2017introduction} reviewed the g-methods.

\end{HomeworkProblem}
